\renewcommand{\teacherTag}{QUESTION 2 \textbullet\ 35 MARKS}
\chapter{The Parking Sensor: Band-pass Sampling and Cross-correlation}

\begin{scenario}
A parking sensor sends out a short 40 kHz ultrasonic pulse, listens for the echo from the nearest obstacle, and works out the distance (Figure~1 of the paper). Because of the Doppler effect and other effects, the received signal occupies the band 38--42 kHz. The Nyquist rate for this signal is therefore 84 kHz, but the sensor samples at a far lower rate, $f_s=18$ kHz. The engineer argues that, in the sampled signal, the copies of the positive and negative halves of the spectrum do not overlap.

Once the echo is captured, its delay is found by cross-correlating the received samples $y[n]$ with the stored samples $x[n]$ of the transmitted pulse,
\[c[m]=\sum_nx[n]\,y[n+m],\]
which has its peak at the delay. The software computes $c[m]$ through the DFT. The delay in samples gives the distance, remembering that the sound travels to the obstacle and back at 340 m/s.
\end{scenario}

\begin{center}
\begin{tikzpicture}[>=Stealth,scale=0.9]
\fill[pattern=north east lines,pattern color=slate](-0.6,-0.05)rectangle(9.3,-0.3);\draw(-0.6,-0.05)--(9.3,-0.05);
\draw[thick,navy,rounded corners=2pt](0,0.35)rectangle(2.6,1.0);\draw[thick,navy](0.6,1.0)--(0.9,1.5)--(1.9,1.5)--(2.2,1.0);
\draw[thick,navy](0.6,0.3)circle(0.28);\draw[thick,navy](2.0,0.3)circle(0.28);
\fill[crimson](2.62,0.62)circle(0.07);\node[font=\tiny,below]at(2.75,0.05){sensor};
\fill[slate!40](8.4,-0.05)rectangle(8.9,2.0);\draw[thick](8.4,-0.05)rectangle(8.9,2.0);\node[font=\scriptsize,above]at(8.65,2.0){obstacle};
\draw[->,thick,teal](2.8,1.15)--node[above,font=\scriptsize]{transmitted pulse $x[n]$}(8.3,1.15);
\draw[<-,thick,coral,dashed](2.8,0.65)--node[below,font=\scriptsize]{echo $y[n]$}(8.3,0.65);
\draw[<->](2.62,-0.6)--node[below,font=\scriptsize]{distance $d$}(8.4,-0.6);
\end{tikzpicture}
\end{center}

%------------------------------------------------------------------ 2(a)
\begin{qcard}{Q2(a)}{Is the engineer right? (band-pass sampling)}{12}
\textbf{Examine}, with a sketch of the spectrum of the received signal sampled at $f_s=18$ kHz, whether the engineer's claim is correct and where the 38--42 kHz band appears in the baseband $\abs f<9$ kHz.
\end{qcard}
\begin{solution}
\step{Step 1 --- The idea.} The Nyquist rate $2f_{\max}=84$ kHz is needed for a signal that occupies \emph{all} frequencies up to 42 kHz. This signal occupies only a 4 kHz-wide band, $B=42-38=4$ kHz. Sampling still repeats the spectrum every $f_s$; what matters is only whether the shifted copies of the band \emph{collide}.

\step{Step 2 --- Where the copies go.} The real signal has two bands: $P=[38,42]$ kHz and its mirror $N=[-42,-38]$ kHz. After sampling, each is copied to every shift by $18k$ kHz:
\begin{center}\small\renewcommand{\arraystretch}{1.25}\begin{tabular}{@{}c|cc@{}}\toprule
shift&copy of $P=[38,42]$&copy of $N=[-42,-38]$\\\midrule
$-18$ / $+18$&$[20,24]$&$[-24,-20]$\\
$-36$ / $+36$&$\mathbf{[2,6]}$&$\mathbf{[-6,-2]}$\\
$-54$ / $+54$&$[-16,-12]$&$[12,16]$\\\bottomrule\end{tabular}\end{center}
Inside the baseband $\abs f<9$ kHz we find exactly two copies: $[2,6]$ kHz (from the positive band) and $[-6,-2]$ kHz (from the negative band). They do not touch: the gap between them is $4$ kHz around $f=0$, and $3$ kHz to each edge $\pm9$.

\begin{center}
\begin{tikzpicture}[>=Stealth]
\begin{scope}[x=0.165cm,y=1cm]
\node[font=\sffamily\scriptsize\bfseries,navy,anchor=west]at(-48,1.4){(i) Received signal (analogue)};
\draw[->](-48,0)--(49,0)node[right,font=\scriptsize]{$f$ (kHz)};
\fill[navy!30](38,0)--(38,0.45)--(42,0.95)--(42,0)--cycle;\draw[navy,thick](38,0)--(38,0.45)--(42,0.95)--(42,0);
\fill[coral!35](-42,0)--(-42,0.95)--(-38,0.45)--(-38,0)--cycle;\draw[coral,thick](-42,0)--(-42,0.95)--(-38,0.45)--(-38,0);
\foreach \x in {-42,-38,38,42}{\node[below,font=\tiny]at(\x,0){$\x$};}\node[below,font=\tiny]at(0,0){$0$};
\node[font=\tiny,navy]at(40,1.15){$P$};\node[font=\tiny,coral]at(-40,1.15){$N$};
\end{scope}
\begin{scope}[x=0.29cm,y=1cm,yshift=-2.6cm]
\node[font=\sffamily\scriptsize\bfseries,navy,anchor=west]at(-27,1.55){(ii) After sampling at $f_s=18$ kHz};
\fill[gold!25](-9,0)rectangle(9,1.3);\node[font=\tiny,anchor=south]at(0,1.3){baseband $\abs f<9$ kHz};
\draw[->](-27,0)--(28,0)node[right,font=\scriptsize]{$f$ (kHz)};
\foreach \c in {20,2,-16}{\fill[navy!30](\c,0)--(\c,0.45)--(\c+4,0.95)--(\c+4,0)--cycle;\draw[navy,thick](\c,0)--(\c,0.45)--(\c+4,0.95)--(\c+4,0);}
\foreach \c in {-24,-6,12}{\fill[coral!35](\c,0)--(\c,0.95)--(\c+4,0.45)--(\c+4,0)--cycle;\draw[coral,thick](\c,0)--(\c,0.95)--(\c+4,0.45)--(\c+4,0);}
\foreach \x in {-24,-16,-12,-9,-6,-2,2,6,9,12,16,24}{\node[below,font=\tiny]at(\x,0){$\x$};}
\draw[dashed,slate](-9,0)--(-9,1.3)(9,0)--(9,1.3);
\end{scope}
\end{tikzpicture}
\end{center}
{\small Blue: copies of the positive band $P$; orange: copies of the negative band $N$. The sloped tops show the orientation (38 kHz edge low, 42 kHz edge high).}

\step{Step 3 --- Formal check (band-pass sampling theorem).} A band $[f_L,f_H]$ sampled at $f_s$ is alias-free if for some integer $n$
\[\frac{2f_H}{n}\le f_s\le\frac{2f_L}{n-1}.\]
With $f_L=38$, $f_H=42$: for $n=5$, $\ \frac{84}5=16.8\le18\le\frac{76}4=19$ \checkmark. So $f_s=18$ kHz lies inside an allowed window.

\step{Step 4 --- Where the band lands.} $38$ kHz maps to $38-2(18)=2$ kHz and $42$ kHz to $6$ kHz: the band is translated, \emph{not} inverted (an even multiple of $f_s/2$ was subtracted). A Doppler echo at $40$ kHz $\pm\Delta$ appears at $4$ kHz $\pm\Delta$.
\begin{keyidea}Band-pass (under-)sampling: a narrow band far from d.c.\ only needs $f_s>2B$ \emph{and} a placement where the copies of $+$ and $-$ halves interleave without touching.\end{keyidea}
\ans{The engineer is right. At $f_s=18$ kHz the band 38--42 kHz appears in the baseband at 2--6 kHz (positive half) and at $-6$ to $-2$ kHz (negative half); the copies do not overlap ($16.8\le f_s\le19$ kHz is an allowed band-pass window).}
\end{solution}

%------------------------------------------------------------------ 2(b)
\begin{qcard}{Q2(b)}{Cross-correlation through the DFT}{12}
\textbf{Demonstrate} that the circular cross-correlation $c[m]=\sum_{n=0}^{N-1}x[n]\,y[((n+m))_N]$ of two real length-$N$ sequences has the DFT $C[k]=X^*[k]\,Y[k]$, stating how much zero-padding, if any, is needed for this route to give the ordinary (linear) cross-correlation.
\end{qcard}
\begin{solution}
\step{Step 1 --- Take the DFT of $c[m]$.} With $W=W_N=e^{-j2\pi/N}$,
\[C[k]=\sum_{m=0}^{N-1}c[m]W^{km}=\sum_{n=0}^{N-1}x[n]\sum_{m=0}^{N-1}y[((n+m))_N]\,W^{km}.\]
\step{Step 2 --- Change the inner index.} Put $l=((n+m))_N$. As $m$ runs over $0..N-1$, so does $l$, and $m\equiv l-n\pmod N$. Since $W^N=1$, $W^{km}=W^{k(l-n)}=W^{kl}W^{-kn}$:
\[C[k]=\sum_{n=0}^{N-1}x[n]W^{-kn}\;\sum_{l=0}^{N-1}y[l]W^{kl}.\]
\step{Step 3 --- Recognise the two factors.} The second sum is $Y[k]$. For a \emph{real} $x$, $\sum x[n]W^{-kn}=\Big(\sum x[n]W^{kn}\Big)^*=X^*[k]$. Hence
\[\boxed{C[k]=X^*[k]\,Y[k]}\qquad\Longrightarrow\qquad c=\mathrm{IDFT}\{X^*\,Y\}.\]
Cost: two FFTs, one product and one inverse FFT, $O(N\log N)$, against $O(N^2)$ for the direct double sum.

\step{Step 4 --- Zero-padding.} The linear cross-correlation $c_{\text{lin}}[m]=\sum_nx[n]y[n+m]$ of a length-$N_x$ pulse and a length-$N_y$ recording is non-zero for lags
\[-(N_x-1)\le m\le N_y-1\qquad(N_x+N_y-1\text{ values}).\]
A length-$M$ circular correlation equals the linear one \emph{folded} with period $M$: $c_{\text{circ}}[m]=\sum_rc_{\text{lin}}[m+rM]$. The folded copies do not overlap only if $M\ge N_x+N_y-1$. Therefore pad both sequences with zeros to
\[\boxed{M\ge N_x+N_y-1\ \ (=2N-1\text{ for two length-}N\text{ sequences})}\]
Then $c_{\text{circ}}[m]=c_{\text{lin}}[m]$ for $0\le m\le N_y-1$, and the negative lags appear at the end: $c_{\text{circ}}[M-l]=c_{\text{lin}}[-l]$. In practice $M$ is rounded up to a power of 2 for the radix-2 FFT.
\ans{$C[k]=\sum_nx[n]W^{-kn}\sum_ly[l]W^{kl}=X^*[k]Y[k]$. Zero-pad both sequences to at least $N_x+N_y-1$ points ($2N-1$ for equal lengths $N$) so the circular correlation equals the linear one.}
\end{solution}

%------------------------------------------------------------------ 2(c)
\begin{qcard}{Q2(c)}{Peak of $c[m]$ and the distance}{11}
To save memory, the sensor only starts recording 30 samples after the pulse has been sent. For one measurement the stored pulse is $x=\{1,2,0,0\}$ and the recording is $y=\{0,1,2,0\}$. Use either method to find the peak of $c[m]$ and, from it, the distance to the obstacle.
\end{qcard}
\begin{solution}
\step{Method 1 --- direct sum.} $c[m]=\sum_nx[n]y[n+m]$, with only $x[0]=1$, $x[1]=2$ non-zero: $c[m]=y[m]+2y[m+1]$.
\begin{center}\small\renewcommand{\arraystretch}{1.25}\begin{tabular}{@{}c|ccccc@{}}
lag $m$&$-1$&0&1&2&3\\\hline
$y[m]+2y[m+1]$&$0+2(0)$&$0+2(1)$&$1+2(2)$&$2+2(0)$&$0$\\
$c[m]$&0&2&\textbf{5}&2&0\end{tabular}\end{center}
\step{Method 2 --- via the DFT.} $N=4$ is enough: the pulse spans 2 samples and the recording 3, so $2+3-1=4$ (part (b)).
\begin{align*}
X[k]&=\{3,\ 1-2j,\ -1,\ 1+2j\},\\
Y[k]&=\{3,\ -2-j,\ 1,\ -2+j\},\\
C[k]=X^*[k]\,Y[k]&=\{9,\ (1+2j)(-2-j),\ (-1)(1),\ (1-2j)(-2+j)\}=\{9,\ -5j,\ -1,\ 5j\}.
\end{align*}
Inverse DFT, $c[m]=\frac14\sum_kC[k]\,j^{km}$:
\[c[0]=\tfrac{9-5j-1+5j}4=2,\quad c[1]=\tfrac{9+5+1+5}4=5,\quad c[2]=\tfrac{9+5j-1-5j}4=2,\quad c[3]=\tfrac{9-5+1-5}4=0 .\]
Both methods agree: $c=\{2,\mathbf 5,2,0\}$, peak at lag $m^\star=1$ (the recorded $\{1,2\}$ is the pulse shifted by one sample).

\begin{center}
\begin{tikzpicture}\begin{axis}[width=7cm,height=4cm,axis lines=middle,xlabel={$m$},ylabel={$c[m]$},xtick={-1,0,1,2,3},ytick={2,5},ymin=0,ymax=6,xmin=-1.6,xmax=3.6,ticklabel style={font=\scriptsize},label style={font=\scriptsize}]
\addplot+[ycomb,navy,thick,mark=*,mark size=2pt]coordinates{(-1,0)(0,2)(1,5)(2,2)(3,0)};
\node[font=\scriptsize,crimson,anchor=west]at(axis cs:1.15,5.2){peak, $m^\star=1$};
\end{axis}\end{tikzpicture}
\end{center}

\step{Step 3 --- Total delay.} The recording started 30 samples after transmission, so the echo delay is
\[D=30+m^\star=31\text{ samples},\qquad \tau=\frac{D}{f_s}=\frac{31}{18\,000}\text{ s}=1.722\text{ ms}.\]
\step{Step 4 --- Distance.} The sound goes to the obstacle and back: $2d=v\tau$.
\[d=\frac{v\tau}{2}=\frac{340\times31}{2\times18\,000}=0.293\text{ m}.\]
\ans{$c[m]=\{2,5,2,0\}$ ($m=0..3$), peak $c[1]=5$; total delay $31$ samples $=1.722$ ms at $f_s=18$ kHz, so the obstacle is $d\approx0.293$ m $\approx29$ cm away.}
\end{solution}
